\begin{table}[htbp]
\centering
\caption{问题2.2：修订后全局枚举与容量约束归属算法复杂度}
\begin{tabular}{|c|c|c|}
\hline
步骤 & 操作 & 复杂度说明 \\
\hline
1. 布局枚举 & 10个小区各取不建、小、中、大四种状态 & $O(4^{10})$ \\
2. 预算剪枝 & 逐布局计算建设成本 & $O(4^{10}\times 10)$ \\
3. 候选归属构造 & 对预算可行布局建立1000米内单站候选 & $O(m\times 10\times p)$ \\
4. 归属优化 & 固定S2估计下在并列最高满意度候选中搜索弹性容量可行归属 & 最坏 $O(m\times p^{10})$ \\
5. 反馈迭代 & 用有效服务负荷更新利用率与响应满意度 & $O(m\times k\times p^{10})$ \\
6. 评价排序 & 计算覆盖、满意度、利润并排序 & $O(n\log n)$ \\
\hline
\end{tabular}

\vspace{0.2cm}
说明：$m$为预算可行布局数，$p$为单个小区可达站点候选数，$k$为方案B归属反馈迭代轮数，$n$为容量可行方案数。本题只有10个小区，归属搜索可在预算剪枝后完成。
\end{table}
